\section{Esoteric Writing}

\begin{quotex}
The prerogative of the human state is objectivity; the essential content of which is the Absolute. There is no knowledge without objectivity of the intelligence; there is no freedom without objectivity of the will; and there is no nobility without objectivity of the soul . . . Esoterism seeks to realize pure and direct objectvity; this is its raison d'être. \flright{\textsc{Frithjof Schuon}}

The vulgar rationalist, even if he is the man most ignorant of all philosophy, is on the contrary the most eager to proclaim himself such at the same time that he dresses himself proudly with the rather ironic title of ``free thinker", while he is in reality merely the slave of all the prejudices current in his time. \flright{\textsc{Rene Guenon}}

Before some audiences not even the possession of the exactest knowledge will make it easy for what we say to produce conviction. For argument based on knowledge implies instruction, and there are people one cannot instruct. \flright{\textsc{Aristotle}}

Education is the only answer to the always pressing question, to the political question par excellence, of how to reconcile order which is not oppression with freedom which is not license. \flright{\textsc{Leo Strauss}}

An exoteric book contains then two teachings: a popular teaching of an edifying character, which is in the foreground and a philosophic teaching concerning the most important subject, which is indicated only between the lines. \flright{\textsc{Leo Strauss}}

The truth is that the modern spirit implies in all who are affected by it in any degree a real hatred of what is secret. \flright{\textsc{Rene Guenon}}

\end{quotex}
A recently published book, \emph{Philosophy Between the Lines: The Lost History of Esoteric Writing}\footnote{\url{https://ndpr.nd.edu/reviews/philosophy-between-the-lines-the-lost-history-of-esoteric-writing/}}, by Arthur Melzer, has brought the ideas of Leo Strauss back to discussion. In the 90s, Strauss was at the center of a small intellectual revival to oppose liberalism and modernity, such as in \emph{Revolt Against Modernity} by \textbf{Ted McAllister}, but Strauss' adventitious association with George Bush's policies discredited him and his theories. The result is that ``conservative" thought in the USA deteriorated into an anti-intellectualism that, on the one hand, worshiped the ideals of Ayn Rand or Horatio Alger, and, on the other, associated the fate of Israel with normative Christianity.

Nevertheless, Strauss, from a secular perspective, did point out the esoteric nature of much philosophical writing. If anyone should doubt it, Melzer has put online an instructive collection of quotes\footnote{\url{http://press.uchicago.edu/sites/melzer/melzer_appendix.pdf}} documenting the extent of esoteric writing. Recently, a Persian expert in the medical theories of Avicenna told me that his texts had borders that ran along the top and the left of the page. The main text on the page said one thing, but the real message was crowded into those borders.

There is resistance to this idea, based on the fear perhaps that some horrendous truth may be unveiled. However, a perusal of the appendix will show instead that the esoteric meaning behind a text is at best simply inaccessible to most people. At worst, they are subject to harmful misinterpretations.

Ultimately, the proof of esoteric writing is learned, not by reading, but by writing. A good writer will seem to say one thing, but mean another. There will be clues and allusions in the text that are often difficult to grasp as a whole. The Germans claim, ``who says A must say B", but often the writer will say A and leave it up to the reader to infer B. Irony may hide the real message.

\paragraph{The Necessity}
The epigraphs have a purpose. The ideal society is a just society, but that requires rulers who are totally objective. How can this be communicated?

People, such as they are, are ruled either by sensation, emotion, or thought. For the first, the only proof required is ``I like it" or ``it feels good." Most people today are ruled by emotions, the most disordered part. What they feel ``with passion" seems true to them. Now, man is a ``rational animal". That is not a definition, nor is rationality something ``added on". Rather, it is what distinguishes men from animals, which live by sensation and emotion, while lacking intellectuality. Therefore, people to a large extent live as animals. It is difficult to control wild animals. They may respond reliably one day, then suddenly change.

Hence, lies are necessary to protect society from erupting. They act as buffers from the clear light of truth.

Even soi-disant free thinkers, as Guenon points out, are subject to the rule of the mass. For them, rationality falls short of pure intellectuality. Thought is merely a tool for partisanship rather than the support for objectivity.

Hence, as Aristotle says, some people are simply not interested in being educated.

\paragraph{Myths and Legends}
A misunderstood aspect of esoteric writing is the idea of the ``noble lie" described by Socrates in the \emph{Republic}. The key adjective here is ``noble". This noble lie is intended for noble purposes of promoting the common good, and the veneration and protection of the city. The effects of denying this ``noble lie" are painfully obvious today, but those details need to wait for another opportunity.

Rather than the term ``noble lie", \textbf{Julius Evola} prefers to call them myths and legends. He accepts Bachofen's characterization:

\begin{quotex}
Symbols and legends, if only in a dramatized form, represent actually and truly the history of the beginnings of a nation, but not the history of events occurring materially on earth, but rather of spiritual processes that have given birth to a new people alongside other people although different in culture and civilization: history, so to say, of its prenatal period. 

\end{quotex}
Those rationalists who can see no deeper than material historical processes will miss this ``third dimension" of history. Myths and legends may be ``falsehoods" in a strictly material sense, but in their spiritual effects, they are much more powerful, hence ``truer".

\paragraph{Esoteric Groups}
So who is eligible to participate in an esoteric group? There can be no \emph{a priori} set of rules, since the presumption is that such men will understand the best rule. We get one clue from Plato. Those who are capable of truly intellectual, objective, and dispassionate thinking will gather together in a special friendship. These friends need the resources and time to participate. Since they are outside the bounds of established institutions and governments, there is the possibility of unintentional conflicts, hence the need for esoteric writing.

At a later period, the Knights Templar took a different approach, although there is not space here to discuss them.

I will give just one more example, the one described by Valentin Tomberg. These ``unknown friends" agree to participate in the stream of Hermetism. By following certain meditative practices, they can reach a certain spiritual depth. Those of a like depth will recognize each other, even if achieved in a different way. Rather than being unnecessarily rebellious, Tomberg instead advocated a close relation to the exoteric institution. This is not unlike what Rene Guenon himself advocated.

There are common myths behind such movements which we may address. Meanwhile, readers can entertain themselves by looking for them first.



\flrightit{Posted on 2014-11-30 by Cologero }

\begin{center}* * *\end{center}

\begin{footnotesize}\begin{sffamily}



\texttt{Tom Walker on 2014-12-01 at 10:53 said: }

But it seems that a lot of Melzer `s referencs are to do with atheists who were intent on hiding their true beliefs e.g Spinoza , Hume etc.


\hfill

\texttt{White Rabbit on 2014-12-01 at 17:23 said: }

Perhaps Melzer references all those atheists as a mere mask for the real esoteric meat of his compilation: a subtle pointing to Traditionalism. The plot thickens!


\hfill

\texttt{Wayne Ferguson on 2014-12-01 at 18:59 said: }

Great article! It's still unclear to me whether he is quoting ``Bachofen" (as you seem to indicate above) OR Cesaro (as you seem to indicate elsewhere). See:

\url{http://www.gornahoor.net/?p=3832}

Also, the phrase, ``history not of living turned materially on earth" seems rather opaque compared to ``the history not of events brought about materially on the earth". Would it be inaccurate to combine the two translations as follows:

``Symbols and legends, if only in a dramatized form, represent actually and truly the history of the beginnings of a nation, but the history not of events brought about materially on the earth, but rather of spiritual processes that have brought to birth, amidst other peoples, a new and different people through culture and civilization: history, so to speak, of its prenatal period."

Thanks…


\hfill

\texttt{Cologero on 2014-12-02 at 23:26 said: }

Good reconstruction of the quote … I updated the post.


\hfill

\texttt{White Rabbit on 2015-02-04 at 23:12 said: }

``Myths and legends may be `falsehoods' in a strictly material sense, but in their spiritual effects, they are much more powerful, hence `truer'."

I for one wonder why that needs to be ``forbidden knowledge."


\hfill

\texttt{White Rabbit on 2015-02-15 at 23:40 said: }

Perhaps it's because faith vivifies the imagination. Or maybe I've introduced more confusion.


\hfill

\texttt{Daniel on 2017-12-02 at 04:00 said: }

It's interesting to read Rene Girard and Leo Strauss against (or alongside) each other on this topic. 

Girard's theory is that sacrificial myths are `lies' told by the lynch mob to mask (even from themselves) the reality of their persecution of a scape goat.

Strauss on the other hand focuses on the need of great thinkers and writers to write `between the lines' for two audiences (one more conventional the other more discerning) to avoid being persecuted for revealing socially unsettling truths.


\end{sffamily}\end{footnotesize}
